% =====================================================================
%  MEMORIA · TRABAJO FINAL DE MÓDULO · PROPUESTA F1
%  Archivo único. Compilar con pdfLaTeX (dos veces para el índice).
%  Sustituye todo el texto que aparece entre [corchetes].
%  Diseño basado en la web personal: IBM Plex Sans / IBM Plex Mono,
%  negro #08080A y azul #E5484D.
% =====================================================================
\documentclass[11pt,a4paper,oneside]{report}

% ---------------- Modo de color ----------------
% \oscurofalse: portada oscura e interior claro (recomendado para imprimir)
% \oscurotrue : todas las páginas con fondo oscuro, como la web
\newif\ifoscuro
\oscurofalse

% ---------------- Idioma y tipografía ----------------
\usepackage[T1]{fontenc}
\usepackage[utf8]{inputenc}
\usepackage[spanish,es-tabla]{babel}
\addto\captionsspanish{\renewcommand{\tablename}{Tabla}\renewcommand{\contentsname}{Índice}}
% IBM Plex Sans (texto y títulos) e IBM Plex Mono (etiquetas y números).
% Vienen con TeX Live completo y Overleaf; si no están, se usan Helvetica y Courier.
\IfFileExists{plex-sans.sty}
  {\usepackage[sfdefault]{plex-sans}}
  {\usepackage[scaled=0.92]{helvet}\renewcommand{\familydefault}{\sfdefault}}
\IfFileExists{plex-mono.sty}
  {\usepackage[scale=0.92]{plex-mono}}
  {\renewcommand{\ttdefault}{pcr}}
\usepackage{microtype}   % \textls{...}: espaciado entre letras de las etiquetas

% ---------------- Página ----------------
\usepackage[a4paper,top=2.8cm,bottom=2.6cm,left=2.8cm,right=2.8cm,headheight=18pt,headsep=14pt]{geometry}
\usepackage{setspace}
\setstretch{1.18}
\setlength{\parindent}{0pt}
\setlength{\parskip}{6pt plus 2pt}

% ---------------- Colores ----------------
\usepackage{xcolor}
% Paleta de la web (fija, se usa en la portada)
\definecolor{negro}{HTML}{08080A}       % fondo
\definecolor{puntos}{HTML}{18181D}      % retícula de puntos
\definecolor{superficie}{HTML}{101014}  % etiquetas y tarjetas
\definecolor{lineaweb}{HTML}{1D1D23}    % líneas finas
\definecolor{bordeweb}{HTML}{2E2E36}    % bordes
\definecolor{blanco}{HTML}{F2F2F4}      % texto principal
\definecolor{grisclaro}{HTML}{9B9BA4}   % segunda línea del título
\definecolor{grismedio}{HTML}{8F8F97}   % texto secundario
\definecolor{azul}{HTML}{488CE5}        % acento sobre fondo oscuro (textos y líneas)
\definecolor{azulboton}{HTML}{3C7DD3}   % botones
\definecolor{azuloscuro}{HTML}{2163BA}  % acento sobre fondo claro (más contraste)
% Colores de las páginas interiores según el modo
\ifoscuro
  \colorlet{fondo}{negro}
  \definecolor{tinta}{HTML}{E4E4E8}
  \colorlet{suave}{grisclaro}
  \colorlet{guia}{grismedio}
  \colorlet{acento}{azul}
  \colorlet{borde}{bordeweb}
  \colorlet{zona}{superficie}
  \colorlet{filaalterna}{superficie}
  \colorlet{cabecera}{lineaweb}
  \colorlet{cabtexto}{blanco}
\else
  \colorlet{fondo}{white}
  \definecolor{tinta}{HTML}{16161A}
  \definecolor{suave}{HTML}{5E5E64}
  \definecolor{guia}{HTML}{74747C}
  \colorlet{acento}{azuloscuro}
  \definecolor{borde}{HTML}{DCDCE1}
  \definecolor{zona}{HTML}{F4F4F6}
  \colorlet{filaalterna}{zona}
  \colorlet{cabecera}{negro}
  \colorlet{cabtexto}{white}
\fi
\color{tinta}            % color por defecto del texto
\ifoscuro\pagecolor{fondo}\fi

% Número con dos cifras (01, 02...); las letras de los anexos se dejan igual
\newcommand{\padnum}[1]{\ifnum0<0#1\relax\ifnum#1<10 0\fi\fi#1}
% Etiqueta como la de la web: línea azul + texto monoespaciado espaciado
\newcommand{\etiqueta}[1]{{\color{acento}\rule[0.55ex]{1.1cm}{0.7pt}\hspace{0.35cm}%
  \ttfamily\small\textls[220]{\MakeUppercase{#1}}}}
% Esquina en escuadra (la de arriba a la izquierda de la web)
\newcommand{\esquina}{\hspace*{-0.55cm}\tikz{\draw[acento, line width=0.8pt]
  (0,0) -- (0,0.42) -- (0.42,0.42);}}
% Encabezado de capítulo: esquina + etiqueta
\newcommand{\rotulocap}[1]{\esquina\par\vspace{6pt}\etiqueta{#1}}

% ---------------- Títulos ----------------
\usepackage{titlesec}
\titleformat{\chapter}[display]
  {\normalfont\sffamily\color{tinta}}
  {\rotulocap{\chaptertitlename\ \padnum{\thechapter}}}
  {14pt}
  {\fontsize{28}{33}\selectfont\bfseries\raggedright}
  [\vspace{10pt}{\color{borde}\titlerule[0.6pt]}]
\titleformat{name=\chapter,numberless}[display]
  {\normalfont\sffamily\color{tinta}}
  {\rotulocap{Trabajo Final de Módulo}}
  {14pt}
  {\fontsize{28}{33}\selectfont\bfseries\raggedright}
  [\vspace{10pt}{\color{borde}\titlerule[0.6pt]}]
\titlespacing*{\chapter}{0pt}{0pt}{22pt}
\titleformat{\section}{\normalfont\sffamily\Large\bfseries\color{tinta}\raggedright}
  {\ttfamily\mdseries\color{acento}\thesection}{0.8em}{}
\titleformat{\subsection}{\normalfont\sffamily\large\bfseries\color{tinta}\raggedright}
  {\ttfamily\mdseries\color{acento}\thesubsection}{0.8em}{}
\titlespacing*{\section}{0pt}{20pt plus 3pt}{6pt}
\titlespacing*{\subsection}{0pt}{14pt plus 2pt}{4pt}
\setcounter{secnumdepth}{2}
\setcounter{tocdepth}{2}

% ---------------- Índice ----------------
\usepackage{titletoc}
\titlecontents{chapter}[2.6em]{\addvspace{10pt}\sffamily\bfseries\color{tinta}}
  {\llap{\makebox[2.6em][l]{\ttfamily\mdseries\color{acento}\padnum{\thecontentslabel}}}}
  {}
  {\hfill\ttfamily\mdseries\color{suave}\thecontentspage}
\titlecontents{section}[5.6em]{\addvspace{1pt}\color{tinta}}
  {\llap{\makebox[3em][l]{\ttfamily\small\color{acento}\thecontentslabel}}}
  {}
  {\titlerule*[0.7pc]{\textcolor{borde}{.}}\ttfamily\small\color{suave}\thecontentspage}
\titlecontents{subsection}[8.8em]{\small\color{tinta}}
  {\llap{\makebox[3.2em][l]{\ttfamily\footnotesize\color{acento}\thecontentslabel}}}
  {}
  {\titlerule*[0.7pc]{\textcolor{borde}{.}}\ttfamily\footnotesize\color{suave}\thecontentspage}

% ---------------- Cabecera y pie ----------------
\usepackage{fancyhdr}
\pagestyle{fancy}
\fancyhf{}
\fancyhead[L]{{\ttfamily\scriptsize\color{suave}\textls[120]{\MakeUppercase{\leftmark}}}}
\fancyhead[R]{\ttfamily\scriptsize\color{acento}\textls[120]{SOC · OLEOSIERRA}}
% Pie: número de página centrado entre dos líneas azules
\newcommand{\barrapie}{{\color{acento}\rule[0.55ex]{0.8cm}{0.6pt}}}
\newcommand{\pie}{%
  \fancyfoot[C]{\barrapie\hspace{8pt}{\ttfamily\footnotesize\color{tinta}\padnum{\thepage}}\hspace{8pt}\barrapie}}
\pie
\renewcommand{\headrulewidth}{0.5pt}
\renewcommand{\headrule}{\ifdim\headrulewidth>0pt{\color{borde}\hrule height\headrulewidth width\headwidth}\fi}
\renewcommand{\chaptermark}[1]{\markboth{\thechapter.\ #1}{}}
\fancypagestyle{plain}{\fancyhf{}\pie\renewcommand{\headrulewidth}{0pt}}

% ---------------- Figuras, tablas y listas ----------------
\usepackage{graphicx}
\usepackage{float}   % [H]: la tabla se queda donde se escribe
\usepackage{caption}
\DeclareCaptionFont{etiq}{\ttfamily\color{acento}}
\DeclareCaptionLabelFormat{mayus}{\textls[120]{\MakeUppercase{#1}~#2}}
\captionsetup{font=small, labelfont={small,etiq}, labelformat=mayus, labelsep=quad}
\usepackage{enumitem}
\setlist{itemsep=2pt,topsep=4pt}
\setlist[itemize]{label=\textcolor{acento}{\rule[0.32ex]{4pt}{4pt}}}
\setlist[enumerate]{label=\textcolor{acento}{\ttfamily\small\padnum{\arabic*}}}
% Tablas con el contenido SIEMPRE centrado (horizontal y vertical):
%   \begin{tabularx}{\linewidth}{C C C} ... \end{tabularx}
\usepackage{array,tabularx,booktabs}
\renewcommand{\tabularxcolumn}[1]{m{#1}}
\newcolumntype{C}{>{\centering\arraybackslash}X}
\newcolumntype{K}{>{\hsize=0.62\hsize\centering\arraybackslash}X}   % columna estrecha centrada
\newcolumntype{Y}{>{\hsize=1.38\hsize\centering\arraybackslash}X}   % columna ancha centrada
\newcolumntype{Q}{>{\hsize=0.75\hsize\centering\arraybackslash}X}   % tabla de 3 columnas: estrecha
\newcolumntype{V}{>{\hsize=0.85\hsize\centering\arraybackslash}X}   % media
\newcolumntype{Z}{>{\hsize=1.40\hsize\centering\arraybackslash}X}   % ancha
% Columna centrada de anchura relativa: P{0.8}, P{1.4}...
% (la suma de los factores debe ser igual al número de columnas)
\newcolumntype{P}[1]{>{\hsize=#1\hsize\centering\arraybackslash}X}
\usepackage{colortbl}
\renewcommand{\arraystretch}{1.35}

% Estilo común de las tablas: \estilotabla (o \estilotabla[1.45] para tablas densas)
\newcommand{\estilotabla}[1][1.7]{%
  \setlength{\aboverulesep}{0pt}\setlength{\belowrulesep}{0pt}%
  \renewcommand{\arraystretch}{#1}}
% Línea superior (azul) e inferior de las tablas
\newcommand{\reglasup}{\noalign{{\color{acento}\hrule height 1.4pt}}}
\newcommand{\reglainf}{\noalign{{\color{tinta}\hrule height 0.8pt}}}
% Celda de cabecera: texto monoespaciado en mayúsculas sobre fondo negro
\newcommand{\cab}[1]{\textcolor{cabtexto}{\ttfamily\footnotesize\bfseries\textls[80]{\MakeUppercase{#1}}}}

% ---------------- Diagrama ----------------
\usepackage{tikz}
\usetikzlibrary{positioning,fit,backgrounds,arrows.meta,calc,babel}

% ---------------- Enlaces ----------------
\usepackage[hidelinks]{hyperref}
\usepackage{xurl}   % permite cortar las URL largas

% ---------------- Huecos para imágenes ----------------
% \imagen{archivo}{pie de figura}{etiqueta}
% Guarda la imagen junto a main.tex (o en una carpeta, p. ej. capturas/proxmox.png).
% Mientras no exista, aparece un recuadro con el nombre del archivo que falta.
\newcommand{\imagen}[3]{%
  \begin{figure}[H]
    \centering
    \IfFileExists{#1}%
      {\includegraphics[width=0.9\linewidth,height=0.42\textheight,keepaspectratio]{#1}}%
      {\begin{tikzpicture}
         \node[draw=borde, dashed, line width=0.6pt, fill=zona,
               minimum width=0.9\linewidth, minimum height=5.2cm, align=center] (h)
           {{\ttfamily\small\bfseries\color{acento}\textls[220]{INSERTAR IMAGEN}}\\[5pt]
            {\ttfamily\scriptsize\color{guia}\detokenize{#1}}};
         \draw[acento, line width=0.8pt] ($(h.north west)+(-4pt,-0.45cm)$) -- ++(0,0.45cm+4pt) -- ++(0.45cm,0);
         \draw[acento, line width=0.8pt] ($(h.south east)+(4pt,0.45cm)$) -- ++(0,-0.45cm-4pt) -- ++(-0.45cm,0);
       \end{tikzpicture}}%
    \caption{#2}\label{#3}%
  \end{figure}}

% Texto guía entre corchetes (en gris para que se vea que hay que sustituirlo)
\newcommand{\guia}[1]{{\color{guia}\itshape[#1]}}

% =====================================================================
\begin{document}

% ============================ PORTADA ================================
% Datos de la portada: cámbialos aquí
\newcommand{\alumno}{Daniel Damota Maldonado}
\newcommand{\profesor}{Carlos Basulto Pardo}
\newcommand{\asignatura}{Incidentes de Ciberseguridad}
\newcommand{\fechaentrega}{[dd/mm/aaaa]}

\begin{titlepage}
\thispagestyle{empty}
\begin{tikzpicture}[remember picture, overlay,
    dato/.style={anchor=north west, inner sep=0pt, text width=7.2cm, align=left}]
\begin{scope}[shift={(current page.north west)}]
  % ---- Fondo y retícula de puntos ----
  \fill[negro] (-0.1,0.1) rectangle (21.1,-29.8);
  \foreach \x in {0.5,1.0,...,20.5}
    \foreach \y in {-0.5,-1.0,...,-29.5}
      \fill[puntos] (\x,\y) circle (0.4pt);
  % Arcos de fondo a la izquierda
  \draw[lineaweb, line width=0.5pt, dash pattern=on 1pt off 2pt] (-4,-15) circle (8.5);
  \draw[lineaweb, line width=0.5pt, dash pattern=on 1pt off 2pt] (-4,-15) circle (11);
  \draw[lineaweb, line width=0.5pt] (21.5,-2.5) circle (6.5);

  % ---- Redes de nodos (fondo) ----
  % Nodo pequeño y nodo con halo
  \tikzset{
    nodo/.pic={\fill[azul] (0,0) circle (1.5pt);},
    nodohalo/.pic={\fill[azul, opacity=0.16] (0,0) circle (8pt);
                   \draw[azul, line width=0.6pt, opacity=0.8] (0,0) circle (4pt);
                   \fill[azul] (0,0) circle (2pt);}}
  % Zona superior derecha
  \coordinate (r1) at (8.0,-2.9);  \coordinate (r2) at (11.0,-0.7);
  \coordinate (r3) at (13.4,-2.4); \coordinate (r4) at (16.9,-1.4);
  \coordinate (r5) at (19.6,-3.1); \coordinate (r6) at (20.7,-0.5);
  \coordinate (r7) at (18.5,-5.7); \coordinate (r8) at (20.3,-8.6);
  \draw[azul, line width=0.5pt, opacity=0.4]
    (r1) -- (r3) -- (r2)  (r3) -- (r4) -- (r5) -- (r6) -- (r4)
    (r5) -- (r7) -- (r8);
  \foreach \p in {r1,r2,r3,r5,r6,r7} \pic at (\p) {nodo};
  \foreach \p in {r4,r8} \pic at (\p) {nodohalo};
  % Franja inferior (entre el texto y los datos)
  \coordinate (s1) at (0.8,-19.0);  \coordinate (s2) at (4.2,-20.9);
  \coordinate (s3) at (7.9,-19.5);  \coordinate (s4) at (11.5,-21.4);
  \coordinate (s5) at (14.6,-19.1); \coordinate (s6) at (17.6,-21.0);
  \coordinate (s7) at (20.5,-18.4);
  \draw[azul, line width=0.5pt, opacity=0.4]
    (s1) -- (s2) -- (s3) -- (s4) -- (s5) -- (s6) -- (s7)
    (s2) -- (s4)  (s3) -- (s5);
  \foreach \p in {s2,s3,s5,s6} \pic at (\p) {nodo};
  \foreach \p in {s1,s4,s7} \pic at (\p) {nodohalo};

  % ---- Esquinas en escuadra ----
  \draw[azul, line width=0.7pt] (1.8,-3.9) -- (1.8,-3.3) -- (2.4,-3.3);
  \draw[azul, line width=0.7pt] (19.2,-19.2) -- (19.2,-19.8) -- (18.6,-19.8);

  % ---- Etiqueta y título ----
  \draw[azul, line width=0.7pt] (2.4,-4.2) -- (3.4,-4.2);
  \node[anchor=west, font=\ttfamily\small, text=azul] at (3.7,-4.2)
    {\textls[300]{TRABAJO FINAL DE MÓDULO}};
  \node[anchor=north west, inner sep=0pt, text width=16.4cm, align=left,
        font=\sffamily\bfseries\fontsize{38}{43}\selectfont, text=blanco] at (2.3,-4.95)
    {\textls[-20]{Centro de Operaciones}\\\textls[-20]{de Seguridad}\\{\color{azul}\textls[-20]{asistido por IA local}}};

  % ---- Párrafo ----
  \node[anchor=north west, inner sep=0pt, text width=8.4cm, align=left,
        font=\sffamily\fontsize{10.5}{16}\selectfont, text=grismedio] (parrafo) at (2.4,-11.6)
    {Propuesta técnica de implantación para \textcolor{blanco}{\bfseries Oleosierra, S.~Coop.~And.}:
     un SOC que detecta los ataques en sus primeras fases, prioriza las alertas con un modelo
     de lenguaje que se ejecuta en la propia cooperativa y documenta cada incidente según la
     normativa vigente.};

  % ---- Botones (siempre 0,9 cm por debajo del párrafo) ----
  \node[anchor=north west, fill=azulboton, inner xsep=12pt, inner ysep=9pt,
        font=\sffamily\small\bfseries, text=white] (b1) at ([yshift=-0.9cm]parrafo.south west) {Propuesta técnica};
  \node[anchor=north west, draw=bordeweb, line width=0.6pt, fill=negro, inner xsep=12pt, inner ysep=9pt,
        font=\sffamily\small\bfseries, text=blanco] at ([xshift=0.3cm]b1.north east)
    {\tikz[baseline=-0.6ex]{\fill[azul, rounded corners=0.8pt] (0,-0.1) rectangle (0.3,0.1);}\hspace{6pt}Oleosierra, S.~Coop.~And.};

  % ---- Retícula circular (desplazada a la derecha y reducida para separarla del texto) ----
% ---- Emblema SOC · IA (Radar SIEM, escudo y malla neuronal) ----
  \begin{scope}[shift={(16.1,-13.9)}, scale=0.82]
    \coordinate (c) at (0,0);

    % Fondo de radar y órbitas de telemetría
    \shade[inner color=azul!28!negro, outer color=superficie] (c) circle (2.85);
    \draw[bordeweb, line width=0.5pt, dash pattern=on 0.8pt off 1.8pt] (c) circle (3.65);
    \draw[azul, line width=0.45pt, opacity=0.6] (c) circle (3.2);
    \draw[azul, line width=1.5pt] ([shift={(115:3.2)}]c) arc[start angle=115, end angle=215, radius=3.2];
    \draw[azul, line width=1.5pt] ([shift={(-65:3.2)}]c) arc[start angle=-65, end angle=35, radius=3.2];

    % Ejes del sonar
    \begin{scope}[azul, line width=0.65pt, opacity=0.85]
      \draw (0,2.85) -- (0,3.95);  \draw (0,-2.85) -- (0,-3.95);
      \draw (-2.85,0) -- (-3.95,0); \draw (2.85,0) -- (3.95,0);
    \end{scope}

    % Escudo exterior
    \draw[fill=superficie, draw=azul, line width=1.05pt, opacity=0.95]
      (0,2.15) -- (1.75,1.45) -- (1.75,-0.2)
      .. controls (1.75,-1.35) and (0.85,-2.05) .. (0,-2.4)
      .. controls (-0.85,-2.05) and (-1.75,-1.35) .. (-1.75,-0.2)
      -- (-1.75,1.45) -- cycle;

    % Hexágono neuronal interno (procesamiento del LLM)
    \begin{scope}[shift={(0,0.1)}]
      \foreach \a in {30,90,...,330} {
        \draw[azul, line width=0.55pt, opacity=0.8] (0,0) -- (\a:1.15);
        \draw[grisclaro, line width=0.4pt, opacity=0.45] (\a:1.15) -- (\a+60:1.15);
      }
      \foreach \a in {30,90,...,330}
        \fill[blanco] (\a:1.15) circle (2.3pt);

      \fill[azul, opacity=0.25] (0,0) circle (9.5pt);
      \fill[azul] (0,0) circle (4.2pt);
      \fill[blanco] (0,0) circle (1.6pt);
    \end{scope}

    % Etiquetas técnicas
    \node[font=\ttfamily\tiny\bfseries, text=azul] at (0,1.62) {\textls[120]{SIEM · IDS}};
    \node[fill=negro, draw=bordeweb, line width=0.5pt, rounded corners=1.5pt,
          inner xsep=5pt, inner ysep=2pt, font=\ttfamily\tiny\bfseries, text=blanco]
          at (0,-1.65) {\textls[140]{SOC · IA}};
  \end{scope}

  % ---- Datos del trabajo ----
  \draw[lineaweb, line width=0.6pt] (2.4,-22.4) -- (18.6,-22.4);
  \node[dato] at (2.4,-23.05)
    {{\ttfamily\scriptsize\color{azul}\textls[220]{ALUMNO}}\\[5pt]{\sffamily\color{blanco}\alumno}};
  \node[dato] at (10.8,-23.05)
    {{\ttfamily\scriptsize\color{azul}\textls[220]{PROFESOR}}\\[5pt]{\sffamily\color{blanco}\profesor}};
  \node[dato] at (2.4,-24.85)
    {{\ttfamily\scriptsize\color{azul}\textls[220]{MÓDULO}}\\[5pt]{\sffamily\color{blanco}\asignatura}};
  \node[dato] at (10.8,-24.85)
    {{\ttfamily\scriptsize\color{azul}\textls[220]{FECHA DE ENTREGA}}\\[5pt]{\sffamily\color{blanco}\fechaentrega}};

  % ---- Pie ----
  \draw[lineaweb, line width=0.6pt] (-0.1,-27.9) -- (21.1,-27.9);
  \node[anchor=west, font=\ttfamily\tiny, text=grismedio] at (2.4,-28.55)
    {\textls[200]{CURSO DE ESPECIALIZACIÓN EN CIBERSEGURIDAD · 2026--2027}};
  \node[anchor=east, font=\ttfamily\tiny, text=grismedio] at (18.6,-28.55)
    {\textls[200]{EMPRESA Y DATOS FICTICIOS}};
\end{scope}
\end{tikzpicture}
\end{titlepage}

% ============================ ÍNDICE =================================
% Páginas del índice sin número (ni cabecera)
\pagenumbering{roman}
\begingroup
  \pagestyle{empty}
  \fancypagestyle{plain}{\fancyhf{}\renewcommand{\headrulewidth}{0pt}}
  % Título del índice sin la etiqueta «Trabajo Final de Módulo»
  \titleformat{name=\chapter,numberless}[block]
    {\normalfont\sffamily\color{tinta}}
    {}
    {0pt}
    {\fontsize{28}{33}\selectfont\bfseries\raggedright}
    [\vspace{10pt}{\color{borde}\titlerule[0.6pt]}]
  \tableofcontents
  \clearpage
\endgroup
\pagenumbering{arabic}

% ============================ 1 =====================================
\chapter{Introducción y justificación}

\section{Contexto de la organización}

Oleosierra, S.~Coop.~And.\ es una cooperativa oleícola \emph{ficticia}, creada para este trabajo, con sede en la comarca de Sierra Mágina (Jaén). Fundada en 1987, agrupa a unos 1.200 socios olivareros, que entregan su aceituna en la almazara de la cooperativa. Allí se molturan entre 18 y 22 millones de kilos de aceituna por campaña, de los que se obtiene aceite de oliva virgen extra que se vende a granel a grandes envasadoras y, desde 2021, también envasado con marca propia a través de una tienda en línea.

La actividad es muy estacional. La campaña de recolección, de noviembre a febrero, concentra la recepción de aceituna, el pesaje en báscula, la molturación y la mayor parte de las liquidaciones a los socios. En esos meses, la plantilla pasa de 38 trabajadores fijos a cerca de 90 personas con la incorporación de personal de campaña. Una parada de pocos días en plena campaña tendría consecuencias económicas graves, porque la aceituna no puede almacenarse sin perder calidad. La Tabla~\ref{tab:empresa} resume sus datos principales.

\begin{table}[htbp]
  \centering
  \caption{Datos principales de Oleosierra (empresa ficticia)}
  \label{tab:empresa}
  \estilotabla
  \begin{tabularx}{\linewidth}{K Y}
    \reglasup
    \rowcolor{cabecera}
    \cab{Concepto} & \cab{Dato} \\
    \rowcolor{filaalterna}
    \textbf{Forma jurídica} & Sociedad cooperativa andaluza \\
    \textbf{Socios} & 1.200 olivareros \\
    \rowcolor{filaalterna}
    \textbf{Plantilla} & 38 fijos; hasta 90 en campaña \\
    \textbf{Facturación anual} & 14 millones de euros \\
    \rowcolor{filaalterna}
    \textbf{Sedes} & Almazara y oficinas en un mismo recinto \\
    \textbf{Canales de venta} & Granel a envasadoras y tienda en línea \\
    \reglainf
  \end{tabularx}
\end{table}

La informática de la cooperativa ha crecido sin planificación. Un único servidor con Windows Server hace de controlador de dominio y de servidor de ficheros, y aloja el programa de gestión de entradas de aceituna y las carpetas con las liquidaciones de los socios. Unos 25 puestos de trabajo con Windows dan servicio a administración, contabilidad, laboratorio de calidad y báscula. El correo electrónico está contratado en la nube y la tienda en línea, alojada en un proveedor externo. Las copias de seguridad se guardan en un NAS conectado a la misma red. El sistema de control de la almazara, de su fabricante, comparte esa red plana con las oficinas, y la gestoría accede por escritorio remoto para preparar las nóminas.

No existe un departamento de seguridad. Un único técnico informático, que también atiende las incidencias diarias, se apoya en un proveedor externo para las tareas complejas. No hay ningún sistema que vigile lo que ocurre en la red ni un procedimiento escrito para actuar ante un incidente.

La cooperativa trata, además, una cantidad importante de datos personales: datos identificativos y bancarios de los socios y de los kilos que entrega cada uno, nóminas y datos de la plantilla (incluido el personal de campaña) y datos de los clientes de la tienda en línea.

\section{Problema y motivación}

Las pymes del sector agroalimentario son un objetivo cada vez más frecuente. Disponen de liquidez y tienen poca tolerancia a las paradas, pero su madurez en ciberseguridad suele ser baja. El panorama de amenazas de ENISA sitúa el ransomware y el phishing entre las amenazas más relevantes para las organizaciones europeas (European Union Agency for Cybersecurity [ENISA], 2025), y ambos encajan con las debilidades de Oleosierra: una red plana, un acceso remoto expuesto, copias de seguridad alcanzables desde la red y personal sin formación específica.

En la última campaña, la cooperativa sufrió dos avisos. La administración recibió un correo que suplantaba a un proveedor de envases y pedía cambiar la cuenta bancaria para el pago de una factura; el fraude se detectó por casualidad, al llamar al proveedor. Semanas después, otra cooperativa de la comarca tuvo que parar su almazara durante varios días por un ataque de ransomware. A raíz de ambos hechos, el consejo rector solicitó una propuesta para mejorar su capacidad de detectar y responder a incidentes.

Se ha elegido la propuesta F1 porque es la que mejor responde a esta necesidad. Un Centro de Operaciones de Seguridad (SOC) permite ver lo que ocurre en la red, detectar ataques en sus primeras fases y actuar con procedimientos definidos, todo ello con el respaldo de un plan de concienciación, un procedimiento de notificación y una estimación de costes que la dirección pueda valorar.

Una cooperativa de este tamaño no puede asumir el coste de un SIEM comercial ni de un SOC externo. Por eso la propuesta se basa en Wazuh, una plataforma de código abierto que reúne en un solo producto la recogida de eventos, la detección basada en reglas, la monitorización de integridad de ficheros, la evaluación de la configuración, la detección de vulnerabilidades y la respuesta activa.

Detectar no es suficiente: también hay que interpretar. Un SIEM genera un volumen de alertas que provoca la llamada \emph{fatiga de alertas}, porque un único técnico no tiene tiempo de estudiarlas todas. Por eso la propuesta incorpora un modelo de lenguaje que resume y prioriza las alertas, explica qué hace un comando sospechoso y propone acciones. El modelo se ejecuta \emph{en local} con Ollama, en un servidor de la propia cooperativa, para que los registros, que contienen datos personales, no salgan de la empresa. La IA solo recomienda y la decisión final siempre corresponde a una persona.

\section{Objetivos}

\textbf{Objetivo general.} Diseñar, implantar en un laboratorio y validar un Centro de Operaciones de Seguridad para Oleosierra basado en Wazuh y asistido por un modelo de lenguaje local para el análisis de alertas, sobre una infraestructura virtualizada en Proxmox VE y segmentada con OPNsense y Suricata, presentado como propuesta técnica a la dirección de la cooperativa.

\textbf{Objetivos específicos:}
\begin{enumerate}[label=\textcolor{acento}{\ttfamily\small\bfseries OE\arabic*}, leftmargin=3.2em]
  \item Analizar los riesgos y la superficie de ataque de la cooperativa e identificar los tres incidentes más probables.
  \item Desplegar la infraestructura virtual del laboratorio en Proxmox VE, con la red segmentada en cuatro VLAN.
  \item Configurar OPNsense como cortafuegos perimetral e inter-VLAN con una política de mínimo privilegio, y Suricata como sistema de detección de intrusiones (y después de prevención) con las reglas ET Open.
  \item Desplegar Wazuh en un nodo único sobre Docker y centralizar en él los eventos de los agentes de Windows y Linux, la telemetría de Sysmon y los registros de OPNsense y Suricata.
  \item Integrar VirusTotal con la monitorización de integridad de ficheros para detectar ficheros maliciosos y eliminarlos mediante respuesta activa.
  \item Desarrollar una integración entre Wazuh y un modelo de lenguaje local (Ollama) que enriquezca las alertas, y evaluar si reduce el trabajo de triaje sin dejar ataques sin detectar.
  \item Validar la solución con siete escenarios de ataque controlados asociados a MITRE ATT\&CK y ajustar las reglas de detección.
  \item Elaborar los playbooks de respuesta, el plan de concienciación y el procedimiento de notificación a las autoridades.
  \item Estimar el coste de implantación y de operación del SOC.
\end{enumerate}

\section{Alcance y limitaciones}

El trabajo comprende el diseño completo del SOC y su validación en un laboratorio. El laboratorio, montado en un único servidor con Proxmox VE, reproduce de forma simplificada la red de la cooperativa con siete máquinas virtuales repartidas en cuatro VLAN: el cortafuegos OPNsense con Suricata; el SIEM y el servidor de IA en la red del SOC; un puesto con Windows~11 y un servidor con Windows Server 2022, que representan los equipos de oficina y el servidor central de la cooperativa; un servidor Linux con servicios web y SSH; y una máquina Kali Linux aislada que simula a un adversario externo. El manual de implantación paso a paso se entrega como documento independiente.

Quedan fuera del alcance:
\begin{itemize}
  \item el despliegue en producción en una empresa real y la alta disponibilidad o el clúster de Wazuh;
  \item la respuesta automática decidida por la IA (el modelo solo sugiere) y el entrenamiento o ajuste fino (\emph{fine-tuning}) del modelo;
  \item el uso de malware real: solo se emplean ficheros de prueba (EICAR) y técnicas simuladas;
  \item el análisis forense profundo de memoria o disco;
  \item la monitorización del sistema de control industrial de la almazara, que se propone como línea futura;
  \item la vigilancia de los servicios en la nube (correo y tienda en línea).
\end{itemize}

Todas las pruebas ofensivas se realizan exclusivamente en el laboratorio, aislado de otras redes mediante OPNsense, y sobre sistemas bajo control del alumno. A VirusTotal solo se envían los hashes de los ficheros, no su contenido, y los registros que analiza la IA no salen del laboratorio. Tanto la empresa como sus datos son ficticios. Por tanto, los resultados demuestran la viabilidad de la propuesta, pero deberían confirmarse en un entorno real antes de su implantación definitiva.

% ============================ 2 =====================================
\chapter{Desarrollo técnico}

Este capítulo recoge la propuesta técnica para Oleosierra. Parte del análisis de riesgos, define la arquitectura de monitorización, describe su implantación y validación en el laboratorio y la completa con los procedimientos organizativos (playbooks, concienciación y notificación) y con la estimación de costes.

\section{Análisis de riesgos y superficie de ataque}
\subsection{Inventario de activos}
\guia{Enumera los activos de la empresa (controlador de dominio, carpeta de liquidaciones, copias de seguridad, puestos de trabajo, correo, sistema de la almazara) y su criticidad.}
\subsection{Superficie de ataque}
\guia{Describe por dónde podría entrar un atacante: superficie externa (correo, accesos remotos, web), humana (personal y temporeros), proveedores y acceso físico.}
\subsection{Matriz de riesgos}
\guia{Explica la metodología (por ejemplo, MAGERIT simplificada: probabilidad × impacto) y presenta la matriz de riesgos con su nivel (bajo, medio, alto o crítico).}
\imagen{capturas/matriz_riesgos.png}{Matriz de riesgos de Oleosierra (probabilidad × impacto)}{fig:matriz}
\subsection{Incidentes más probables}
\guia{A partir de la matriz, identifica y justifica los tres incidentes más probables en el sector: ransomware, phishing con fraude por correo y acceso no autorizado con credenciales robadas. Indica qué escenarios de prueba de la sección~\ref{sec:validacion} reproducen cada uno.}

% ---------------------------------------------------------------------
\section{Arquitectura de monitorización}
\subsection{Visión general}

La solución se despliega en un único servidor físico con Proxmox VE (Figura~\ref{fig:arquitectura}). OPNsense actúa como única puerta de enlace: separa la red del laboratorio de la red exterior (WAN) y enruta y filtra el tráfico entre cuatro VLAN, que viajan por un \emph{bridge} con soporte de VLAN. Suricata inspecciona el tráfico de las interfaces internas, y todos los eventos (los de los agentes instalados en los equipos y los del propio cortafuegos) convergen en Wazuh. Wazuh consulta la reputación de los ficheros sospechosos en VirusTotal y envía las alertas relevantes al modelo de lenguaje local, cuyo análisis vuelve al panel para que el analista decida.

\begin{figure}[htbp]
\centering
% El diagrama se coloca con posiciones relativas (cada caja se sitúa respecto a la anterior
% y las zonas se ajustan a su contenido) y se escala al ancho de la página: así no se
% solapa aunque cambie la fuente.
\resizebox{\linewidth}{!}{%
\begin{tikzpicture}[
    font=\sffamily\scriptsize,
    equipo/.style={draw=tinta!70, line width=0.7pt, fill=fondo, rounded corners=2pt,
                   text width=3.2cm, minimum height=1.35cm, align=center, inner sep=3pt},
    externo/.style={draw=suave, dashed, line width=0.7pt, fill=fondo, rounded corners=2pt,
                    text width=3.6cm, minimum height=0.9cm, align=center, inner sep=3pt},
    zona/.style={draw=borde, line width=0.6pt, fill=zona, rounded corners=4pt, inner sep=7pt},
    rotulo/.style={font=\sffamily\scriptsize\bfseries, text=acento, align=center,
                   text width=3.2cm, inner sep=1pt},
    enlace/.style={draw=tinta, line width=0.6pt},
    flujo/.style={draw=acento, line width=0.8pt, dashed, -{Stealth[length=2.2mm]}},
    nota/.style={font=\ttfamily\tiny, text=acento, inner sep=2pt}
  ]
  % ---- Columnas de las VLAN (separadas 4,3 cm) ----
  \node[rotulo] (r10) at (0,0)    {VLAN 10 · SOC\\[1pt]{\ttfamily\mdseries\tiny 10.10.10.0/24}};
  \node[rotulo] (r20) at (4.3,0)  {VLAN 20 · Endpoints\\[1pt]{\ttfamily\mdseries\tiny 10.10.20.0/24}};
  \node[rotulo] (r30) at (8.6,0)  {VLAN 30 · Servidores\\[1pt]{\ttfamily\mdseries\tiny 10.10.30.0/24}};
  \node[rotulo] (r99) at (12.9,0) {VLAN 99 · Atacante\\[1pt]{\ttfamily\mdseries\tiny 10.10.99.0/24 · aislada}};

  % ---- Equipos (cada uno debajo de su rótulo o del equipo anterior) ----
  \node[equipo, below=4mm of r10]    (wazuh)  {\textbf{wazuh-aio} · .10\\[1pt]Wazuh 4.14 (Docker)\\[1pt]{\tiny manager · indexer · panel}};
  \node[equipo, below=8mm of wazuh]  (ollama) {\textbf{ollama} · .20\\[1pt]Ollama + LLM 7--8B\\[1pt]{\tiny API REST :11434}};
  \node[equipo, below=4mm of r20]    (win11)  {\textbf{win11} · .11\\[1pt]Windows 11 Pro\\[1pt]{\tiny Sysmon · FIM $\rightarrow$ VirusTotal}};
  \node[equipo, below=8mm of win11]  (winsrv) {\textbf{winsrv} · .12\\[1pt]Windows Server 2022\\[1pt]{\tiny Sysmon · eventos 4625}};
  \node[equipo, below=4mm of r30]    (web)    {\textbf{ubuntu-web} · .10\\[1pt]Ubuntu Server 24.04\\[1pt]{\tiny web + SSH · respuesta activa}};
  \node[equipo, below=4mm of r99]    (kali)   {\textbf{kali} · .10\\[1pt]Kali Linux\\[1pt]{\tiny nmap · hydra}};

  % ---- Zonas: se ajustan automáticamente a su contenido ----
  \begin{scope}[on background layer]
    \node[zona, fit=(r10)(wazuh)(ollama)] (z10) {};
    \node[zona, fit=(r20)(win11)(winsrv)] (z20) {};
    \node[zona, fit=(r30)(web)]           (z30) {};
    \node[zona, fit=(r99)(kali)]          (z99) {};
  \end{scope}

  \draw[flujo, {Stealth[length=2.2mm]}-{Stealth[length=2.2mm]}] (wazuh) -- node[nota, right, xshift=2pt] {análisis} (ollama);
  \draw[draw=acento, line width=0.8pt, -{Stealth[length=2.2mm]}] (kali.west) -- (web.east);

  % ---- Troncal 802.1Q ----
  \coordinate (tr) at (0,1.45);
  \draw[tinta, line width=2pt] ($(z10.north west |- tr)+(0.3,0)$) -- ($(z99.north east |- tr)+(-0.3,0)$);
  \foreach \z in {z10,z20,z30,z99} \draw[enlace] (\z.north) -- (\z.north |- tr);
  \node[nota, anchor=south east] at ($(z99.north east |- tr)+(-0.3,0)$) {vmbr1 · trunk 802.1Q (VLAN 10/20/30/99)};

  % ---- Cortafuegos ----
  \node[equipo, fill=acento, text=white, text width=4.8cm, minimum height=1.0cm] (fw) at (6.45,2.75)
    {\textbf{OPNsense 26.x + Suricata}\\[1pt]{\tiny Cortafuegos · IDS/IPS · ET Open}};
  \draw[enlace] (fw.south) -- (fw.south |- tr);

  % ---- Exterior ----
  \node[externo] (inet) at (6.45,4.75) {Internet / red del centro\\[1pt]{\tiny WAN · vmbr0}};
  \node[externo] (vt)   at (12.1,4.75) {\textbf{VirusTotal API v3}\\[1pt]{\tiny reputación de hashes}};
  \draw[enlace] (inet) -- (fw);
  \draw[flujo] (fw.east) -| node[nota, pos=0, anchor=south west, xshift=3pt] {hashes · 443/TCP} (vt.south);

  % ---- Leyenda de flujos (debajo de las VLAN 30 y 99) ----
  \node[draw=borde, fill=fondo, rounded corners=2pt, anchor=north, text width=7.9cm,
        align=left, font=\sffamily\tiny, inner sep=5pt] (leyenda)
        at ($(z30.south)!0.5!(z99.south)+(0,-0.4)$)
    {\textbf{\color{acento}Flujos principales (controlados por OPNsense)}\\[3pt]
     Agentes $\rightarrow$ wazuh-aio: 1514/TCP (eventos) y 1515/TCP (registro)\\[1pt]
     OPNsense y Suricata $\rightarrow$ wazuh-aio: syslog 514/UDP\\[1pt]
     wazuh-aio $\rightarrow$ ollama: 11434/TCP (solo desde 10.10.10.10)\\[1pt]
     wazuh-aio $\rightarrow$ VirusTotal: 443/TCP (solo hashes)\\[1pt]
     Kali $\rightarrow$ VLAN 20 y 30: tráfico de ataque inspeccionado por Suricata};

  % ---- Host Proxmox: rodea todo lo que se ejecuta en el servidor ----
  \begin{scope}[on background layer]
    \node[draw=suave, line width=0.8pt, rounded corners=5pt, inner sep=10pt,
          fit=(fw)(z10)(z20)(z30)(z99)(leyenda)] (host) {};
  \end{scope}
  \node[anchor=north west, font=\sffamily\scriptsize\bfseries, inner sep=6pt] at (host.north west)
    {Host físico · Proxmox VE 9.x};
\end{tikzpicture}}
\caption{Arquitectura lógica del laboratorio y flujos principales}
\label{fig:arquitectura}
\end{figure}

\subsection{Segmentación de la red}

Cada VLAN representa una zona de la red de la cooperativa (Tabla~\ref{tab:vlan}). En una implantación real, esta segmentación sustituiría a la red plana actual de Oleosierra y separaría también el sistema de control de la almazara en una VLAN propia, que queda fuera del alcance del laboratorio.

\begin{table}[H]
  \centering
  \caption{VLAN del laboratorio y su equivalencia en la cooperativa}
  \label{tab:vlan}
  \small
  \estilotabla
  \begin{tabularx}{\linewidth}{P{0.75} P{0.85} P{1.4}}
    \reglasup
    \rowcolor{cabecera}
    \cab{VLAN} & \cab{Red} & \cab{Equipos y función} \\
    \rowcolor{filaalterna}
    \textbf{10 · SOC} & 10.10.10.0/24 & SIEM (wazuh-aio) y servidor de IA (ollama) \\
    \textbf{20 · Endpoints Windows} & 10.10.20.0/24 & Puesto de usuario y servidor Windows: red de oficina \\
    \rowcolor{filaalterna}
    \textbf{30 · Servidores} & 10.10.30.0/24 & Servidor Linux con servicios web y SSH \\
    \textbf{99 · Atacante} & 10.10.99.0/24 & Kali Linux, aislada del SOC \\
    \reglainf
  \end{tabularx}
\end{table}

OPNsense aplica una política de \textbf{denegación por defecto}: solo se permiten los flujos imprescindibles y se registra todo lo demás (Tabla~\ref{tab:politica}). La red del SOC es la más protegida: el modelo de lenguaje solo es accesible desde el SIEM, y la máquina atacante no puede alcanzar ningún equipo de monitorización.

\begin{table}[H]
  \centering
  \caption{Política de filtrado entre VLAN en OPNsense}
  \label{tab:politica}
  \footnotesize
  \estilotabla[1.5]
  \begin{tabularx}{\linewidth}{P{1.05} P{0.85} P{0.95} P{0.75} P{1.4}}
    \reglasup
    \rowcolor{cabecera}
    \cab{Origen} & \cab{Destino} & \cab{Puerto / servicio} & \cab{Acción} & \cab{Justificación} \\
    \rowcolor{filaalterna}
    VLAN 20 y 30 (agentes) & wazuh-aio & 1514 y 1515/TCP & Permitir & Envío de eventos y registro de agentes \\
    OPNsense & wazuh-aio & 514/UDP (syslog) & Permitir & Registros del cortafuegos y alertas de Suricata \\
    \rowcolor{filaalterna}
    Equipo de administración (VLAN 10) & wazuh-aio & 443/TCP & Permitir & Acceso al panel de Wazuh \\
    wazuh-aio & ollama & 11434/TCP & Permitir & API del modelo; inaccesible desde otras VLAN \\
    \rowcolor{filaalterna}
    wazuh-aio & Internet (VirusTotal) & 443/TCP & Permitir & Consulta de hashes y actualizaciones \\
    VLAN 99 (Kali) & VLAN 20 y 30 & Servicios expuestos & Permitir (solo en pruebas) & Escenarios de ataque inspeccionados por Suricata \\
    \rowcolor{filaalterna}
    VLAN 99 (Kali) & VLAN 10 (SOC) & Cualquiera & Denegar & Proteger la infraestructura de monitorización \\
    Cualquier VLAN & Cualquier destino & Resto & Denegar y registrar & Mínimo privilegio y trazabilidad \\
    \reglainf
  \end{tabularx}
\end{table}

\subsection{Componentes tecnológicos}

Todos los componentes son gratuitos o de código abierto, lo que elimina el coste de licencias, uno de los principales frenos para que una pyme disponga de un SOC. La Tabla~\ref{tab:componentes} recoge la función de cada uno.

\begin{table}[H]
  \centering
  \caption{Componentes tecnológicos del SOC}
  \label{tab:componentes}
  \footnotesize
  \estilotabla[1.5]
  \begin{tabularx}{\linewidth}{P{0.7} P{0.8} P{1.5}}
    \reglasup
    \rowcolor{cabecera}
    \cab{Componente} & \cab{Versión} & \cab{Función en el proyecto} \\
    \rowcolor{filaalterna}
    \textbf{Proxmox VE} & 9.x & Hipervisor de tipo 1, con bridges con soporte de VLAN y snapshots para restaurar el laboratorio antes de cada escenario \\
    \textbf{OPNsense} & 26.x (rama estable) & Cortafuegos, enrutamiento entre VLAN, DHCP y DNS; envía sus registros a Wazuh por syslog \\
    \rowcolor{filaalterna}
    \textbf{Suricata} & Integrado en OPNsense & Detección de intrusiones en red con reglas ET Open: primero como IDS y después como IPS \\
    \textbf{Wazuh} & 4.14.x (Docker, nodo único) & SIEM/XDR: decodificadores y reglas, FIM, evaluación de configuración, detección de vulnerabilidades, MITRE ATT\&CK y respuesta activa \\
    \rowcolor{filaalterna}
    \textbf{Docker Engine y Compose} & Versión estable & Despliegue reproducible del manager, el indexer y el dashboard con volúmenes persistentes \\
    \textbf{Sysmon} & Versión vigente & Telemetría avanzada de Windows: procesos, conexiones de red, registro y acceso a procesos sensibles \\
    \rowcolor{filaalterna}
    \textbf{VirusTotal API v3} & API pública & Reputación de los hashes de los ficheros detectados por la monitorización de integridad \\
    \textbf{Ollama y LLM} & Modelo de 7--8B cuantizado & Análisis y explicación de alertas en local mediante su API REST \\
    \rowcolor{filaalterna}
    \textbf{Kali Linux y Atomic Red Team} & Versiones vigentes & Generación controlada de actividad maliciosa alineada con MITRE ATT\&CK \\
    \reglainf
  \end{tabularx}
\end{table}

\subsection{SIEM: Wazuh}

Un SIEM (\emph{Security Information and Event Management}) centraliza los registros de los equipos y servicios, los normaliza y los correlaciona mediante reglas para generar alertas. Se ha elegido Wazuh porque es de código abierto y reúne en una sola plataforma la detección basada en reglas con su técnica de MITRE ATT\&CK asociada, la monitorización de integridad de ficheros (FIM), la evaluación de la configuración (SCA), la detección de vulnerabilidades y la respuesta activa. Splunk se licencia según el volumen de datos y Elastic Security reserva parte de sus funciones a las suscripciones de pago, mientras que Wazuh ofrece todas sus capacidades sin coste de licencias.

Wazuh se despliega en un único nodo sobre Docker Compose, con sus tres componentes en contenedores: \textbf{wazuh.manager}, que recibe y analiza los eventos; \textbf{wazuh.indexer}, que los almacena e indexa; y \textbf{wazuh.dashboard}, la interfaz web, accesible por HTTPS. Los volúmenes persistentes conservan la configuración y los datos al recrear los contenedores, y certificados propios cifran las comunicaciones entre componentes.

Las fuentes de eventos son:
\begin{itemize}
  \item \textbf{Agentes de Windows} (win11 y winsrv): registros de seguridad, como los inicios de sesión fallidos (evento 4625), y la telemetría de Sysmon.
  \item \textbf{Agente de Linux} (ubuntu-web): registros del sistema y de autenticación (\texttt{sshd}).
  \item \textbf{OPNsense y Suricata}: registros del cortafuegos y alertas del IDS, recibidos por syslog.
  \item \textbf{VirusTotal}: resultado de la consulta de reputación de los ficheros nuevos o modificados.
\end{itemize}

\subsection{Telemetría de endpoints: Sysmon}

Los registros nativos de Windows no recogen con detalle lo que ejecuta un atacante. Sysmon, de Microsoft Sysinternals, añade eventos de creación de procesos con su línea de comandos (evento~1), conexiones de red (evento~3), acceso a la memoria de otros procesos, como LSASS (evento~10), y modificaciones del registro (evento~13). Se parte de una configuración comunitaria contrastada (SwiftOnSecurity o sysmon-modular), que filtra el ruido habitual del sistema y conserva los eventos útiles para la detección.

\subsection{Detección de intrusiones: Suricata}

Suricata, integrado en OPNsense, inspecciona el tráfico de las interfaces internas, es decir, el que circula entre las VLAN. Utiliza el conjunto de reglas \textbf{ET Open} de Emerging Threats, que detecta escaneos, ataques a servicios, tráfico de malware conocido y comunicaciones con infraestructuras maliciosas, y se actualiza de forma programada desde el propio cortafuegos. Sus alertas llegan a Wazuh por syslog.

Suricata se activa primero en modo \textbf{IDS}, que solo alerta, para identificar y ajustar las reglas que generan falsos positivos. Una vez ajustadas, pasa a modo \textbf{IPS}, en el que bloquea en línea el tráfico malicioso. Así se evita que el sistema de prevención corte tráfico legítimo, algo especialmente grave en plena campaña.

\subsection{Inteligencia de amenazas}

La propuesta combina dos fuentes de inteligencia de amenazas:
\begin{itemize}
  \item \textbf{ET Open}, en la red: las firmas de Suricata incorporan indicadores de amenazas conocidas.
  \item \textbf{VirusTotal}, en los equipos: cuando la monitorización de integridad detecta un fichero nuevo o modificado en un directorio de riesgo, Wazuh envía su hash a VirusTotal. Si varios motores antivirus lo identifican como malicioso, se genera una alerta y la respuesta activa elimina el fichero.
\end{itemize}

La API pública de VirusTotal limita las consultas a 4 por minuto y 500 al día, por lo que solo se monitorizan los directorios de mayor riesgo (Descargas y carpetas temporales). Además, sus condiciones de uso no permiten emplearla en servicios comerciales, de modo que en la implantación real en la cooperativa habría que contratar la API de pago o sustituirla por una fuente gratuita que lo permita. Esta decisión se tiene en cuenta en la estimación de costes.

% ---------------------------------------------------------------------
\section{Implantación en laboratorio}
\subsection{Metodología}
\guia{Describe las fases de la implantación y el hito de verificación de cada una.}
\subsection{Máquinas virtuales y direccionamiento}
\guia{Tabla con las máquinas virtuales, su sistema operativo, dirección IP y recursos. Requisitos del host.}
\subsection{Proxmox VE y red virtual}
\guia{Instalación del host, bridges con soporte de VLAN y snapshots.}
\subsection{OPNsense y Suricata}
\guia{Interfaces VLAN, DHCP y DNS, reglas del cortafuegos, pruebas de segmentación y configuración de Suricata.}
\subsection{Wazuh en Docker}
\guia{Preparación del host, certificados, Docker Compose y acceso al panel.}
\subsection{Agentes y fuentes de eventos}
\guia{Agentes de Windows y Linux, Sysmon, syslog de OPNsense y decodificadores de Suricata.}
\subsection{VirusTotal, FIM y respuesta activa}
\guia{Directorios monitorizados, integración con la API, respuesta activa y prueba con EICAR.}
\subsection{Reglas de detección propias}
\guia{Reglas y decodificadores propios con su técnica de MITRE ATT\&CK.}

% ---------------------------------------------------------------------
\section{Triaje de alertas asistido por IA}
\subsection{Funcionamiento}
\guia{Describe la integración entre Wazuh y Ollama: qué alertas se envían al modelo, cómo se construye el prompt y cómo vuelve el análisis al panel. Incluye un diagrama del flujo.}
\subsection{Decisiones de diseño y salvaguardas}
\guia{Disparo selectivo, fallo seguro, aislamiento, protección frente a la inyección de instrucciones y supervisión humana.}
\subsection{Selección del modelo}
\guia{Modelos comparados, criterios y modelo elegido.}

% ---------------------------------------------------------------------
\section{Validación con escenarios de ataque}
\label{sec:validacion}
\subsection{Escenarios}
\guia{Tabla con los escenarios de prueba (técnica de MITRE ATT\&CK, herramienta, fuente de detección y resultado esperado) y relación con los incidentes más probables.}
\subsection{Resultados}
\guia{Criterios de éxito fijados en el anteproyecto y resultados obtenidos, con un análisis crítico.}

% ---------------------------------------------------------------------
\section{Playbooks de respuesta a incidentes}
\subsection{Organización de la respuesta}
\guia{Define los niveles de escalado (L1, L2 y L3), quién interviene en cada uno y la matriz RACI de responsabilidades. Sigue las fases de NIST SP 800-61r3.}
\subsection{Playbook 1: ransomware}
\guia{Detalla cada fase (detección, análisis, contención, erradicación, recuperación, notificación y lecciones aprendidas) con sus acciones, responsables y herramientas. Indica qué alertas del SOC lo activan (escenarios E5 a E7) y qué acciones de contención ofrece Wazuh.}
\imagen{capturas/playbook_ransomware.png}{Diagrama de flujo del playbook de ransomware}{fig:pb1}
\subsection{Playbook 2: phishing y fraude por correo}
\guia{Misma estructura que el playbook anterior, adaptada a este incidente (escenario E4).}
\imagen{capturas/playbook_phishing.png}{Diagrama de flujo del playbook de phishing y fraude por correo}{fig:pb2}
\subsection{Playbook 3: acceso no autorizado}
\guia{Misma estructura que el playbook anterior, adaptada a este incidente (escenarios E2 y E3).}
\imagen{capturas/playbook_acceso.png}{Diagrama de flujo del playbook de acceso no autorizado}{fig:pb3}
\subsection{Simulacro}
\guia{Documenta al menos un simulacro de mesa: fecha, participantes, escenario, tiempos, decisiones tomadas y mejoras detectadas.}

\section{Plan de concienciación}
\guia{Identifica los colectivos de la empresa y su riesgo principal, define las acciones formativas con su calendario (sesiones, píldoras, simulacro de phishing) y los indicadores para medir su eficacia (tasa de clics, correos reportados, plantilla formada). Incluye la formación del técnico en el uso del analista de IA (alfabetización en IA, capítulo~\ref{cap:normativa}).}
\imagen{capturas/simulacro_phishing.png}{Correo utilizado en el simulacro de phishing}{fig:simulacro}

\section{Procedimiento de notificación a autoridades}
\guia{Explica cómo se clasifica un incidente (taxonomía y niveles de la Guía nacional de notificación y gestión de ciberincidentes) y a quién se notifica, en qué plazo y por qué canal: AEPD (72 horas), personas afectadas, INCIBE-CERT, CSIRT de referencia si aplica NIS2, Fuerzas y Cuerpos de Seguridad y aseguradora. Describe también el comité de crisis y la comunicación interna.}
\imagen{capturas/flujo_notificacion.png}{Flujo de notificación y comité de crisis}{fig:notificacion}

\section{Estimación del coste de implantación}
\guia{Calcula la inversión inicial: servidor con los requisitos del laboratorio (8 núcleos y 16 hilos, 64 GB de RAM, 1 TB NVMe y GPU opcional para Ollama), horas de implantación y formación. El software es gratuito, pero en producción hay que contar con la API de pago de VirusTotal (o una alternativa con uso comercial permitido). Calcula también el coste anual de operación (horas del técnico, mantenimiento y consumo eléctrico) y compáralo con contratar un SOC gestionado y con el coste de sufrir un incidente en plena campaña.}

% ============================ 3 =====================================
\chapter{Marco normativo aplicable}
\label{cap:normativa}

Este capítulo analiza las normas que Oleosierra debe cumplir, qué le exigen y por qué le afectan, además de la norma ISO/IEC 27001, cuya certificación facilita el SOC. La situación normativa corresponde a la fecha de redacción (octubre de 2026). La Tabla~\ref{tab:normativa} resume el resultado al final del capítulo.

\section{Protección de datos: RGPD y LOPDGDD}

\textbf{¿Se aplica? Sí, de forma obligatoria.} El Reglamento General de Protección de Datos (RGPD) y la Ley Orgánica 3/2018 (LOPDGDD) se aplican a cualquier organización que trate datos personales. Oleosierra es responsable del tratamiento de tres grandes conjuntos de datos: los de sus 1.200 socios (identificación, cuenta bancaria, kilos entregados y liquidaciones), los de su plantilla, incluido el personal de campaña (nóminas, contratos, datos de salud laboral), y los de los clientes de la tienda en línea.

\textbf{Qué le exige y cómo lo cubre la propuesta:}
\begin{itemize}
  \item \textbf{Integridad y confidencialidad} (art.~5.1.f) y \textbf{seguridad del tratamiento} (art.~32): aplicar medidas técnicas y organizativas adecuadas al riesgo, entre ellas la capacidad de detectar y restaurar el servicio tras un incidente. El SOC, los playbooks y la estrategia de copias de seguridad responden directamente a esta obligación.
  \item \textbf{Responsabilidad proactiva} (art.~24): poder demostrar que se cumple. El registro centralizado de eventos del SIEM y la documentación de cada incidente sirven como evidencia.
  \item \textbf{Notificación de brechas a la AEPD} (art.~33): sin dilación indebida y, a más tardar, en 72 horas desde que se tiene constancia, salvo que sea improbable que suponga un riesgo. Sin un sistema de detección, la cooperativa podría tardar semanas en saber que ha sufrido una brecha; el SOC reduce ese tiempo y el procedimiento de notificación ordena la respuesta.
  \item \textbf{Comunicación a los afectados} (art.~34): cuando la brecha suponga un alto riesgo para los socios, trabajadores o clientes.
\end{itemize}

El propio SOC trata datos personales, porque los registros incluyen nombres de usuario, direcciones IP y actividad de los empleados. Este tratamiento se basa en el interés legítimo en garantizar la seguridad de la red (art.~6.1.f y considerando~49 del RGPD) y debe respetar el principio de minimización y unos plazos de conservación definidos. Por el mismo principio, el diseño evita comunicar datos a terceros: la integración con VirusTotal solo envía el hash de los ficheros, nunca su contenido, y el modelo de IA se ejecuta en un servidor propio. Además, la LOPDGDD (art.~87) y el Estatuto de los Trabajadores exigen informar a la plantilla de que el uso de los equipos de la empresa puede ser supervisado con fines de seguridad, por lo que la implantación del SOC debe ir acompañada de una política de uso de los dispositivos digitales.

Oleosierra no está obligada a designar un delegado de protección de datos (DPD), porque no figura entre los supuestos del artículo~37 del RGPD ni del artículo~34 de la LOPDGDD. Aun así, se recomienda designar uno externo por el volumen de datos de socios que maneja.

\section{Directiva NIS2}

\textbf{¿Se aplica? Todavía no es exigible en España, pero lo será.} La Directiva (UE) 2022/2555 (NIS2) amplía las obligaciones de ciberseguridad a nuevos sectores. Su anexo~II incluye la \emph{producción, transformación y distribución de alimentos}, en concreto las empresas alimentarias dedicadas a la transformación industrial y a la distribución al por mayor. Oleosierra encaja en ambas actividades: transforma la aceituna en aceite en su almazara y lo vende a granel a envasadoras.

La directiva solo se aplica a medianas y grandes empresas. Con 38 trabajadores fijos y el personal de campaña, Oleosierra suma unas 55 unidades de trabajo anual, y su facturación (14 millones de euros) supera el umbral de 10 millones. Por tanto, es una \textbf{mediana empresa} y, por su sector, sería una \textbf{entidad importante}.

Sin embargo, España no ha completado la transposición. El anteproyecto de Ley de Coordinación y Gobernanza de la Ciberseguridad se aprobó en el Consejo de Ministros el 14 de enero de 2025, pero en octubre de 2026 aún no se ha publicado en el BOE, y la Comisión Europea ha llevado a España ante el Tribunal de Justicia de la UE por el retraso. Mientras tanto sigue vigente la normativa anterior (Real Decreto-ley 12/2018), que no incluye al sector alimentario.

\textbf{Por qué se tiene en cuenta igualmente:} cuando la ley entre en vigor, la cooperativa tendrá que cumplir, entre otras, estas obligaciones:
\begin{itemize}
  \item \textbf{Gobernanza} (art.~20): el consejo rector deberá aprobar las medidas de ciberseguridad, supervisarlas y recibir formación, y podrá responder por su incumplimiento.
  \item \textbf{Medidas de gestión de riesgos} (art.~21): análisis de riesgos, gestión de incidentes, continuidad de negocio, seguridad de la cadena de suministro, formación y control de accesos, entre otras.
  \item \textbf{Notificación de incidentes significativos} (art.~23) al CSIRT de referencia: alerta temprana en 24 horas, notificación en 72 horas e informe final en un mes.
\end{itemize}
Las entidades importantes se enfrentan a multas de hasta 7 millones de euros o el 1,4\,\% de su facturación mundial. La propuesta se diseña desde el principio para cumplir estas obligaciones, de modo que la cooperativa no tenga que improvisar cuando se apruebe la ley.

\section{ISO/IEC 27001:2022}

La norma ISO/IEC 27001 es voluntaria, pero las grandes envasadoras y cadenas de distribución a las que vende Oleosierra exigen cada vez más garantías de seguridad a sus proveedores. Por eso, obtener la certificación sería una ventaja comercial a medio plazo.

El SOC propuesto \textbf{facilita ese proceso de certificación}, porque ya cumple varios de los controles del anexo~A que el auditor revisará:
\begin{itemize}
  \item \textbf{Gestión de incidentes} (A.5.24 a A.5.28): planificación, valoración de eventos, respuesta, lecciones aprendidas y recogida de evidencias, mediante los playbooks y el registro de cada incidente.
  \item \textbf{Inteligencia de amenazas} (A.5.7), mediante las reglas ET Open de Suricata y la consulta de reputación en VirusTotal.
  \item \textbf{Concienciación y formación} (A.6.3) y \textbf{notificación de eventos por parte del personal} (A.6.8), mediante el plan de concienciación.
  \item \textbf{Protección contra el malware} (A.8.7), \textbf{gestión de vulnerabilidades técnicas} (A.8.8) y \textbf{gestión de la configuración} (A.8.9), mediante la integración con VirusTotal, la respuesta activa y los módulos de vulnerabilidades y de evaluación de la configuración de Wazuh.
  \item \textbf{Registro de la actividad} (A.8.15) y \textbf{monitorización} (A.8.16), mediante el SIEM y el sistema de detección de intrusiones.
  \item \textbf{Seguridad de las redes} (A.8.20) y \textbf{segregación de redes} (A.8.22), mediante las VLAN y la política de filtrado de OPNsense.
\end{itemize}
Además, el análisis de riesgos del capítulo~2 sirve de base para la evaluación y el tratamiento de riesgos que exige la norma. Así, la cooperativa tendría ya implantada y documentada una parte importante del sistema de gestión antes de iniciar la certificación.

\section{Reglamento de Inteligencia Artificial}

\textbf{¿Se aplica? Sí, con obligaciones limitadas.} El Reglamento (UE) 2024/1689 regula los sistemas de IA según su nivel de riesgo. Al utilizar el analista de IA, Oleosierra actúa como \emph{responsable del despliegue}.

El analista de IA \textbf{no es un sistema de alto riesgo}: no figura entre los usos del anexo~III (empleo, crédito, educación, biometría, infraestructuras críticas, etc.), porque se limita a ayudar a clasificar alertas de seguridad de una red corporativa y no toma decisiones sobre personas. Tampoco interactúa con el público, por lo que no le afectan las obligaciones de transparencia del artículo~50.

La obligación que sí le afecta es la \textbf{alfabetización en IA} (art.~4), aplicable desde el 2 de febrero de 2025: la empresa debe adoptar medidas para que las personas que usan el sistema entiendan cómo funciona y cuáles son sus limitaciones. El Reglamento (UE) 2026/1744 (ómnibus digital sobre IA) ha suavizado su redacción, que ahora es una obligación de adoptar medidas y no de garantizar un nivel concreto, pero sigue vigente. La propuesta la cubre con la formación del técnico y con la documentación del sistema. Además, aunque no sea obligatorio, se aplica el principio de \textbf{supervisión humana}: la IA solo recomienda y ninguna acción se ejecuta sin que una persona la revise.

\filbreak   % si el resumen no cabe entero, pasa completo a la página siguiente
\section{Resumen}

La Tabla~\ref{tab:normativa} resume qué normas afectan a Oleosierra, si son obligatorias y por qué.

\begin{table}[H]
  \centering
  \caption{Resumen del marco normativo aplicable a Oleosierra}
  \label{tab:normativa}
  \small
  \estilotabla
  \begin{tabularx}{\linewidth}{Q V Z}
    \reglasup
    \rowcolor{cabecera}
    \cab{Norma} & \cab{¿Obligatoria?} & \cab{Motivo} \\
    \rowcolor{filaalterna}
    \textbf{RGPD y LOPDGDD} & Sí & Trata datos de socios, plantilla y clientes \\
    \textbf{NIS2} & Sí, cuando España la transponga & Sector alimentario y mediana empresa: entidad importante \\
    \rowcolor{filaalterna}
    \textbf{ISO/IEC 27001} & No (voluntaria) & El SOC cumple controles del anexo~A y facilita la certificación \\
    \textbf{Reglamento de IA} & Sí, con obligaciones limitadas & Despliega un sistema de IA que no es de alto riesgo \\
    \reglainf
  \end{tabularx}
\end{table}

% ============================ 4 =====================================
\chapter{Conclusiones}
\section{Resultados obtenidos}
\guia{Valora en qué medida se han cumplido los objetivos OE1 a OE9, con los resultados de las pruebas y del resto del trabajo.}
\section{Limitaciones}
\guia{Reflexiona sobre las limitaciones del trabajo: laboratorio frente a empresa real, número de escenarios, ataques conocidos de antemano, tamaño de los modelos y ejecución sin GPU, y límites de la API pública de VirusTotal.}
\section{Líneas de mejora}
\guia{Propón cómo continuar el trabajo: respuesta automática supervisada, asistente conversacional de threat hunting, nuevas fuentes de registros (correo en la nube, NAS de copias), más modelos o vigilancia de la red industrial de la almazara.}

% ===================== DECLARACIÓN DE USO DE IA ======================
\chapter*{Declaración de uso de IA generativa}
\addcontentsline{toc}{chapter}{Declaración de uso de IA generativa}
\markboth{Declaración de uso de IA generativa}{}

\guia{Declara las herramientas de IA generativa utilizadas, en qué partes del trabajo y de qué forma.}

% ============================ 5 =====================================
\chapter{Referencias bibliográficas}

\begin{list}{}{\leftmargin=1.27cm \itemindent=-1.27cm \itemsep=6pt}
  \item Directiva (UE) 2022/2555 del Parlamento Europeo y del Consejo, de 14 de diciembre de 2022, relativa a las medidas destinadas a garantizar un elevado nivel común de ciberseguridad en toda la Unión (Directiva SRI 2). (2022). \textit{Diario Oficial de la Unión Europea}, L 333, 80--152. \url{https://eur-lex.europa.eu/eli/dir/2022/2555/oj}
  \item European Union Agency for Cybersecurity. (2025). \textit{ENISA threat landscape 2025}. \url{https://www.enisa.europa.eu/publications/enisa-threat-landscape-2025}
  \item International Organization for Standardization. (2022). \textit{Information security, cybersecurity and privacy protection --- Information security management systems --- Requirements} (Norma ISO/IEC 27001:2022). \url{https://www.iso.org/standard/27001}
  \item Ley Orgánica 3/2018, de 5 de diciembre, de Protección de Datos Personales y garantía de los derechos digitales. (2018). \textit{Boletín Oficial del Estado}, 294. \url{https://www.boe.es/eli/es/lo/2018/12/05/3}
  \item Microsoft. (s.~f.). \textit{Sysmon}. Microsoft Learn. \url{https://learn.microsoft.com/es-es/sysinternals/downloads/sysmon}
  \item MITRE Corporation. (s.~f.). \textit{MITRE ATT\&CK}. \url{https://attack.mitre.org}
  \item OPNsense. (s.~f.). \textit{OPNsense documentation: Intrusion prevention system}. \url{https://docs.opnsense.org}
  \item Proofpoint. (s.~f.). \textit{Emerging Threats Open ruleset}. \url{https://rules.emergingthreats.net}
  \item Proxmox Server Solutions. (s.~f.). \textit{Proxmox VE administration guide}. \url{https://pve.proxmox.com/pve-docs/}
  \item Reglamento (UE) 2016/679 del Parlamento Europeo y del Consejo, de 27 de abril de 2016, relativo a la protección de las personas físicas en lo que respecta al tratamiento de datos personales y a la libre circulación de estos datos. (2016). \textit{Diario Oficial de la Unión Europea}, L 119, 1--88. \url{https://eur-lex.europa.eu/eli/reg/2016/679/oj}
  \item Reglamento (UE) 2024/1689 del Parlamento Europeo y del Consejo, de 13 de junio de 2024, por el que se establecen normas armonizadas en materia de inteligencia artificial (Reglamento de Inteligencia Artificial). (2024). \textit{Diario Oficial de la Unión Europea}. \url{https://eur-lex.europa.eu/eli/reg/2024/1689/oj}
  \item Reglamento (UE) 2026/1744 del Parlamento Europeo y del Consejo, por el que se modifica el Reglamento (UE) 2024/1689 en lo que respecta a la simplificación de la aplicación de las normas armonizadas en materia de inteligencia artificial (ómnibus digital sobre IA). (2026). \textit{Diario Oficial de la Unión Europea}. \url{https://eur-lex.europa.eu/eli/reg/2026/1744/oj}
  \item VirusTotal. (s.~f.). \textit{VirusTotal API v3 reference}. \url{https://docs.virustotal.com}
  \item Wazuh. (s.~f.). \textit{Wazuh documentation: Deployment on Docker, VirusTotal integration and Sysmon monitoring}. \url{https://documentation.wazuh.com/current/}
\end{list}

\guia{Añade las fuentes adicionales que cites al completar el capítulo 2, en formato APA 7.ª edición y manteniendo el orden alfabético.}

% ============================ ANEXOS ================================
\appendix
\renewcommand{\appendixname}{Anexo}
\chapter{Capturas y evidencias del laboratorio}
\guia{Capturas de pantalla ordenadas por capítulo, cada una con su pie de figura.}

\chapter{Ficheros de configuración}
\guia{Extractos comentados o enlace al repositorio con: docker-compose.yml, bloques relevantes de ossec.conf, reglas y decodificadores propios, script custom-ollama y exportación de las reglas de OPNsense. No incluyas la clave de la API de VirusTotal.}

\chapter{Informes de los incidentes simulados}
\label{anx:informes}
\guia{Un informe por escenario (E1 a E7): detección, clasificación según la Guía nacional de notificación y gestión de ciberincidentes, contención, erradicación, recuperación y lecciones aprendidas, junto con el análisis de la IA y su comparación con el análisis manual.}

\chapter{Playbooks en formato visual}
\guia{Diagramas de flujo de los tres playbooks.}

\chapter{Materiales de concienciación}
\guia{Infografías, guiones de las píldoras formativas y correo del simulacro de phishing.}

\chapter{Plantillas de notificación}
\guia{Formulario de notificación de brechas de la AEPD cumplimentado, aviso a los afectados y comunicación interna.}

\chapter{Relación con los resultados de aprendizaje}
\guia{Tabla que relacione el trabajo con los resultados de aprendizaje RA1 a RA5 del módulo (Real Decreto 479/2020).}


\end{document}
